% AMD DCGPU System Integration notes:
% - Post-silicon Datacenter GPU System Integration team for MI300X / MI325X / MI350X.
% - Used an internal orchestration framework to run workloads across performance, HPC, and ML applications.
% - Validation domains: RAS, resets, IO/ExtIO, performance, system management.
% - Worked across diverse Linux/Windows platforms with varied BMCs, SBIOS, and OS images (including RHEL/CentOS).
% - Participated in hardware and software debugs spanning ROCm, PMFW, firmware, drivers, and platform issues.
% - Built Python tooling to accelerate workload job launch and creation of system-level integration test cases.
% - Used Docker for application environments, GitHub for code, and Jira for tracking work.
% - Hosted meetings, reviewed firmware-related requirements with senior engineers, and led process-improvement initiatives within the team.

% Prior role title (comment out; do not delete): Datacenter GPU Infrastructure \& Validation Engineer (Co-op)
% Prior split BMC/BIOS segment (comment out; do not delete):
% {Advanced Micro Devices (AMD) $|$ \textbf{BMC}, \textbf{BIOS} $|$ \textbf{Python}, \textbf{Linux}, \textbf{Git}, \textbf{Bash}, \textbf{ROCm}, \textbf{Jira}}
% Prior standard heading with Jira in header (comment out; do not delete):
% \resumeSubheading
%   {Advanced Micro Devices (AMD) $|$ \textbf{BMC}, \textbf{BIOS}, \textbf{Python}, \textbf{Linux}, \textbf{Git}, \textbf{Bash}, \textbf{ROCm}, \textbf{Jira}}{Markham, ON}
% Prior smaller keyword header (comment out; do not delete):
% \resumeSubheadingSmall {Advanced Micro Devices (AMD) $|$ ...}{Markham, ON}{...}{...}
% Prior header ordering BMC-first without C/C++, Validation split, or Integration tail (comment out; do not delete):
% {Advanced Micro Devices (AMD) $|$ BMC, BIOS, Python, Linux, Git, Bash, ROCm, Performance, IO, RAS, Memory, Windows, Docker}
% Prior single-argument header all same font size (comment out; do not delete):
% \resumeSubheadingProminent {Advanced Micro Devices (AMD) $|$ Python, C/C++, ...}{Markham, ON}{Co-op title}{dates}
\resumeSubheadingProminent
  {Advanced Micro Devices (AMD)}
  % Prior keyword tail with Git, Docker, ROCm (comment out; do not delete) — NVIDIA targeting prefers leaner pipe:
  % {$|$ Python, C/C++, BMC, OOB, BIOS, Linux, Bash, Git, Windows, ROCm, Docker, RAS, IO, Performance, Memory, System Management, firmware, system integration}
  % Prior keyword tail with OOB and system integration tokens (comment out; do not delete):
  % {$|$ Python, C/C++, BMC, OOB, BIOS, Linux, Bash, Windows, RAS, IO, Performance, Memory, System Management, firmware, system integration}
  % {$|$ Python, C/C++, BMC, BIOS, UEFI, Linux, Bash, Windows, ROCm, RAS, IO, Performance, Memory, System Management, firmware}
  % Prior keyword tail pre-Graphics IP Validation AMD 2026-05-06 (comment out; do not delete):
  % {$|$ Python, C/C++, BMC, BIOS/UEFI, Linux, Bash, Windows, ROCm, RAS, IO, Performance, Memory, System Management}
  % Prior keyword tail AMD Graphics IP Validation 2026-05-06 (comment out; do not delete):
  % {$|$ Python, C/C++, Verilog, Quartus, GPU, post-silicon, BMC, BIOS/UEFI, Linux, Windows, ROCm, RAS, IO, System Management}
  % Prior keyword tail pre-Docker/ROCm/PyTorch header (comment out; do not delete):
  % {$|$ Python, C/C++, GPU, Linux, Bash, Windows, Verilog, Quartus, test tooling, BMC, BIOS/UEFI, RAS, IO}
  % Prior keyword tail NVIDIA GPU Architecture NCG 2026-05-06 (comment out; do not delete):
  % {$|$ Python, C/C++, GPU, Linux, Docker, ROCm, PyTorch, BMC, BIOS/UEFI, RAS, IO}
  % Prior keyword tail NVIDIA Low Power ASIC Engineer NCG 2026-05-06 (comment out; do not delete):
  % {$|$ Python, C/C++, Verilog, Quartus, GPU validation, Linux, Windows, BMC, BIOS/UEFI, Performance, System Management}
  % {$|$ Python, C/C++, Verilog, Quartus, testbenches, post-silicon GPU, Memory, IO, firmware, Linux, BMC, BIOS/UEFI}
  % {$|$ Python, C/C++, ROCm, Linux, GitHub, Docker, post-silicon GPU, profiling workflows, APIs, firmware, BMC, BIOS/UEFI}
  % {$|$ Python, C/C++, ROCm, Linux, GitHub, Docker, post-silicon GPU validation, firmware, BMC, BIOS/UEFI}
  % {$|$ Python, C/C++, Linux, Docker, Git, Bash, Windows, ROCm, BMC, BIOS/UEFI}
  % Prior keyword tail pre Senior Software Systems Validation Engineer GPU Linux 2026-05-20 (comment out; do not delete):
  % {$|$ Python, C/C++, Linux, Docker, Git, Bash, Windows, ROCm, BMC, BIOS/UEFI, RAS, IO, PPA, Memory, System Management}
  % Prior keyword tail pre guidelines AMD bullet rewrite 2026-05-20 (comment out; do not delete):
  % {$|$ Python, C/C++, Linux, CentOS, ROCm, KVM, QEMU, Docker, Git, Bash, Windows, BMC, BIOS/UEFI, RAS, IO, GFX, PPA, Memory, System Management}
  % {$|$ Python, C/C++, Linux, Ubuntu, RHEL, CentOS, BKC, ROCm, KVM, QEMU, Docker, Git, Bash, Windows, BMC, BIOS/UEFI, RAS, IO, GFX, PPA, Memory, System Management}
  {$|$ Python, C/C++, Linux, Windows, Hardware Architecture, Platform Integration, Test Planning, Operational Validation, BMC, BIOS/UEFI, Out-of-Band Management, System Design, Firmware, Docker}
  {Markham, ON}
  {System Integration \& Validation Engineer}
  {May 2024 -- Aug 2025}

\resumeItemListStart

% Prior MAAS-first layout (comment out; do not delete) — Canonical MAAS targeting:
% \resumeSubheading
%   {Advanced Micro Devices (AMD) $|$ \textbf{Python}, \textbf{Linux}, \textbf{MAAS}, \textbf{Bash}, \textbf{Windows}}{Markham, ON}
% \resumeItem{Debugged \textbf{MAAS} provisioning flows for MI300X, MI325X, and MI350X systems, resolving commissioning and imaging issues so GPU validation racks stayed ready for recurring automated integration runs.}

% Prior ops-first wording (comment out; do not delete):
% \resumeItem{Operated automated post-silicon validation on \textbf{Linux} datacenter GPU platforms for MI300X, MI325X, and MI350X against validated BKCs, running integration passes across Performance, IO, and platform quality at 100+ executions per week.}

% Prior Linux-only diagnostic bullet (comment out; do not delete):
% \resumeItem{Ran automated post-silicon diagnostic validation across MI300X, MI325X, and MI350X \textbf{Linux} platforms against validated BKCs, sustaining Performance, IO, and platform reliability coverage at roughly 100+ executions weekly.}

% Prior diagnostic bullet without explicit system integration (comment out; do not delete):
% \resumeItem{Ran automated post-silicon diagnostic validation across MI300X, MI325X, and MI350X \textbf{Linux} and \textbf{Windows} platforms against validated BKCs, sustaining \textbf{Performance}, \textbf{IO}, \textbf{RAS}, \textbf{Memory}, and \textbf{System Management} coverage at roughly 100+ executions weekly.}

% Prior diagnostic bullet overlapping lab bullet on OS and validation pillars (comment out; do not delete):
% \resumeItem{Ran automated post-silicon diagnostic validation across MI300X, MI325X, and MI350X \textbf{Linux} and \textbf{Windows} platforms against validated BKCs, sustaining \textbf{Performance}, \textbf{IO}, \textbf{RAS}, \textbf{Memory}, and \textbf{System Management} coverage while exercising datacenter GPU \textbf{system integration} and \textbf{HPC} workload paths at roughly 100+ executions weekly.}

% Prior denser diagnostic bullet (comment out; do not delete):
% \resumeItem{Ran automated post-silicon diagnostic validation across MI300X, MI325X, and MI350X \textbf{Linux} and \textbf{Windows} platforms against validated BKCs, sustaining \textbf{RAS}, \textbf{IO}, \textbf{Performance}, \textbf{Memory}, and \textbf{System Management} coverage while exercising datacenter GPU \textbf{system integration} and \textbf{HPC} workload paths at roughly 100+ executions weekly.}

% Prior longer diagnostic bullet (comment out; do not delete):
% \resumeItem{Ran automated weekly diagnostics on MI300X, MI325X, and MI350X datacenter GPUs under \textbf{Linux} and \textbf{Windows}, using validated platform baselines to sustain \textbf{RAS}, \textbf{IO}, \textbf{Performance}, \textbf{Memory}, and \textbf{System Management} while running \textbf{system integration} and \textbf{HPC} workloads at roughly 100+ test runs per week.}

% Prior automated-diagnostics style bullet (comment out; do not delete):
% \resumeItem{Ran automated diagnostics on MI300X, MI325X, and MI350X datacenter GPUs on \textbf{Linux} and \textbf{Windows} against validated baselines, sustaining \textbf{RAS}, \textbf{IO}, \textbf{Performance}, \textbf{Memory}, and \textbf{System Management} plus \textbf{system integration} and \textbf{HPC} at 100+ runs weekly.}

% Prior bullet with HPC bold and unbolded validation pillars (comment out; do not delete):
% \resumeItem{Ran \textbf{HPC} and ML stress workloads on MI300X, MI325X, and MI350X under \textbf{Linux} and \textbf{Windows} while test modules exercised RAS, IO, Performance, Memory, and System Management alongside those runs, sustaining validated baseline coverage at 100+ runs weekly.}

% Prior bullet without test-lead orchestration framing (comment out; do not delete):
% \resumeItem{Ran HPC and ML stress workloads on MI300X, MI325X, and MI350X under \textbf{Linux} and \textbf{Windows} while test modules exercised \textbf{RAS}, \textbf{IO}, \textbf{Performance}, \textbf{Memory}, and \textbf{System Management} alongside those runs, sustaining validated baseline coverage at \textbf{100+ runs weekly}.}

% Prior lead bullet without explicit post-silicon GPU validation phrasing (comment out; do not delete):
% \resumeItem{Acted as rotating \textbf{test lead} for MI300X, MI325X, and MI350X lanes, orchestrating \textbf{100+ runs weekly} on \textbf{Linux} and \textbf{Windows} while HPC and ML stress workloads ran and companion test modules exercised \textbf{RAS}, \textbf{IO}, \textbf{Performance}, \textbf{Memory}, and \textbf{System Management} under validated baselines.}

% Prior lead bullet without Docker, PyTorch, and ROCm ML debug phrasing (comment out; do not delete):
% \resumeItem{Acted as rotating \textbf{test lead} for MI300X, MI325X, and MI350X \textbf{post-silicon} datacenter \textbf{GPU} validation lanes, orchestrating \textbf{100+ runs weekly} on \textbf{Linux} and \textbf{Windows} while HPC and ML stress workloads ran and companion test modules exercised \textbf{RAS}, \textbf{IO}, \textbf{Performance}, \textbf{Memory}, and \textbf{System Management} under validated baselines.}

% Prior combined lead bullet with ML debug on same line (comment out; do not delete):
% \resumeItem{Acted as rotating \textbf{test lead} for MI300X, MI325X, and MI350X \textbf{post-silicon} datacenter \textbf{GPU} validation lanes, orchestrating \textbf{100+ runs weekly} on \textbf{Linux} and \textbf{Windows} while HPC and ML stress workloads ran and companion test modules exercised \textbf{RAS}, \textbf{IO}, \textbf{Performance}, \textbf{Memory}, and \textbf{System Management} under validated baselines, debugging ML failures through \textbf{Docker} images, \textbf{PyTorch}, and \textbf{ROCm} during those runs.}

% \resumeItem{Acted as rotating \textbf{test lead} for MI300X, MI325X, and MI350X \textbf{post-silicon} datacenter \textbf{GPU} validation lanes, orchestrating \textbf{100+ runs weekly} on \textbf{Linux} and \textbf{Windows} while HPC and ML stress workloads ran and companion test modules exercised \textbf{RAS}, \textbf{IO}, \textbf{Performance}, \textbf{Memory}, and \textbf{System Management} under validated baselines.}
% \resumeItem{Acted as rotating \textbf{test lead} for MI300X, MI325X, and MI350X post-silicon GPU validation, coordinating \textbf{100+ runs weekly} on \textbf{Linux} and \textbf{Windows} across HPC and ML workload gates.}
% \resumeItem{Acted as rotating \textbf{test lead} for MI300X, MI325X, and MI350X semiconductor post-silicon programs, coordinating \textbf{100+ diagnostic runs weekly} on \textbf{Linux} and \textbf{Windows} test stations across \textbf{RAS}, \textbf{IO}, \textbf{PPA}, \textbf{Memory}, and \textbf{System Management} coverage.}
% \resumeItem{Rotated as \textbf{test lead} on MI300X, MI325X, and MI350X post-silicon lanes; tracked \textbf{100+ diagnostic runs weekly} on \textbf{Linux} and \textbf{Windows} stations across \textbf{RAS}, \textbf{IO}, \textbf{PPA}, \textbf{Memory}, and \textbf{System Management}.}
% Prior GPU-specific validation framing (comment out; do not delete):
% \resumeItem{Rotated as \textbf{test lead} on MI300X, MI325X, and MI350X post-silicon lanes; tracked \textbf{100+ diagnostic runs weekly} on \textbf{Linux} and \textbf{Windows} stations across \textbf{RAS}, \textbf{Int-IO}, \textbf{Ext-IO}, \textbf{GFX}, \textbf{PPA}, \textbf{Memory}, and \textbf{System Management}.}
\resumeItem{Led hardware platform readiness and test execution for new MI300X, MI325X, and MI350X datacenter GPUs, orchestrating cross-functional alignment with infrastructure and validation teams to integrate new platforms into operational environments; sustained \textbf{100+ validation runs weekly} to verify hardware quality and readiness for datacenter deployment.}
% Prior short lead bullet (comment out; do not delete):
% \resumeItem{Led platform readiness and \textbf{test execution} for new hardware integration across MI300X, MI325X, and MI350X systems, orchestrating \textbf{100+ validation runs weekly} on \textbf{Linux} and \textbf{Windows} to verify hardware platform quality for datacenter deployment.}

% Prior ML triage bullet spanning MI300X, MI325X, and MI350X (comment out; do not delete):
% \resumeItem{Triaged and reproduced ML workload failures across \textbf{Docker}-based stacks, \textbf{PyTorch} execution, and the \textbf{ROCm} runtime on MI300X, MI325X, and MI350X GPUs during overlapping \textbf{HPC} and ML stress campaigns, narrowing software versus hardware root cause before reruns.}

% Prior ML stack triage bullet MI350X (NVIDIA GPU Arch NCG; comment out; do not delete):
% \resumeItem{Triaged and reproduced ML workload failures across \textbf{Docker}-based stacks, \textbf{PyTorch} execution, and the \textbf{ROCm} runtime on \textbf{MI350X} systems during overlapping HPC and ML stress campaigns, narrowing software versus hardware root cause before reruns.}

% Prior Performance / System Management stress bullet NVIDIA Low Power ASIC Engineer NCG (comment out; do not delete):
% \resumeItem{Stressed \textbf{Performance} and \textbf{System Management} paths on MI300X, MI325X, and MI350X \textbf{GPU} silicon while HPC and ML workloads ran, surfacing power- and performance-adjacent platform regressions before wider integration sign-off.}

% \resumeItem{Ran post-silicon validation focused on \textbf{Memory} and \textbf{IO} coverage on MI300X, MI325X, and MI350X datacenter \textbf{GPUs} beside HPC and ML workloads, surfacing interface-level regressions before wider integration gates.}
% \resumeItem{Coordinated feature enablement with firmware, validation, and infrastructure partners across MI300X, MI325X, and MI350X profiling and workload paths, keeping cross-team integration milestones aligned for release cycles.}
% Prior ML triage bullet spanning MI300X, MI325X, and MI350X (comment out; do not delete):
% \resumeItem{Triaged and reproduced ML workload failures across \textbf{Docker} images and the \textbf{ROCm} runtime on MI300X, MI325X, and MI350X systems, narrowing software versus hardware root cause before reruns.}
% \resumeItem{Deployed and maintained \textbf{Linux} and \textbf{Windows} GPU lab racks with \textbf{BMC}, \textbf{BIOS/UEFI}, and \textbf{Docker}-based workload images, supporting NPI bring-up and repeatable test-station configuration across datacenter GPU platforms.}
% \resumeItem{Maintained \textbf{Linux} and \textbf{Windows} GPU lab racks with \textbf{CentOS}/RHEL images, \textbf{BMC}, \textbf{BIOS/UEFI} flashes, and \textbf{Docker} workload images for NPI bring-up and repeatable test-station configs on datacenter GPU platforms.}
% \resumeItem{Set up \textbf{KVM/QEMU} on \textbf{MI325X} \textbf{CentOS} hosts, wrote bring-up docs the team reused, and ran validation and debug on guest GPU paths before wider lane sign-off.}
% \resumeItem{Updated \textbf{Linux} GPU racks on \textbf{Ubuntu}, \textbf{RHEL}, and \textbf{CentOS} and \textbf{Windows} systems to validated \textbf{BKC} baselines; flashed \textbf{BMC} and \textbf{BIOS/UEFI}, staged \textbf{Docker} images, and enabled \textbf{NPI} so ML workloads could run on \textbf{MI350X}.}
% \resumeItem{Set up \textbf{KVM/QEMU} on \textbf{MI325X} \textbf{CentOS} hosts and wrote setup docs the team reused for guest bring-up.}
% Prior NPI-specific GPU phrasing (comment out; do not delete):
% \resumeItem{Updated \textbf{Linux} GPU racks on \textbf{Ubuntu}, \textbf{RHEL}, and \textbf{CentOS} and \textbf{Windows} systems to validated \textbf{BKC} baselines; flashed \textbf{BMC} and \textbf{BIOS/UEFI}, staged \textbf{Docker} images, and configured test configs and ML workload parameters so \textbf{NPI}-partitioned workloads could run on \textbf{MI350X}.}
\resumeItem{Developed and documented hardware platform bring-up procedures for \textbf{Linux} and \textbf{Windows} systems, including \textbf{BMC} and \textbf{BIOS/UEFI} configuration, validated baselines, and \textbf{Docker} workload staging; mentored 3+ team members on hardware integration best practices, reducing new-platform onboarding time and improving distributed team consistency.}
\resumeItem{Set up \textbf{KVM/QEMU} on \textbf{MI325X} \textbf{CentOS} hosts, wrote setup docs the team reused for running workloads, validated \textbf{MI325X} KVM guests under production-style workloads, and triaged failures found during those runs.}

% Prior longer lead bullet with workload images and call-path phrasing (comment out; do not delete):
% \resumeItem{Acted as rotating \textbf{test lead} for MI300X, MI325X, and MI350X \textbf{post-silicon} datacenter \textbf{GPU} validation lanes, orchestrating \textbf{100+ runs weekly} on \textbf{Linux} and \textbf{Windows} while HPC and ML stress workloads ran and companion test modules exercised \textbf{RAS}, \textbf{IO}, \textbf{Performance}, \textbf{Memory}, and \textbf{System Management} under validated baselines, debugging ML regressions across \textbf{Docker} workload images, \textbf{PyTorch} call paths, and the \textbf{ROCm} runtime when issues appeared in those runs.}

% Prior lab bullet without system integration phrasing (comment out; do not delete):
% \resumeItem{Deployed and maintained \textbf{Linux} and \textbf{Windows} GPU lab stacks spanning \textbf{BMC}, \textbf{BIOS}, and OS images plus \textbf{Docker}-based workload environments to stress \textbf{RAS}, \textbf{IO}, \textbf{Performance}, \textbf{Memory}, and \textbf{System Management}, trimming early bring-up failures by about 20\%.}

% Prior lab bullet naming Docker (comment out; do not delete):
% \resumeItem{Deployed and maintained \textbf{Linux} and \textbf{Windows} GPU lab stacks spanning \textbf{BMC}, \textbf{BIOS}, and OS images plus \textbf{Docker}-based workload environments for repeatable \textbf{system integration} bring-up, stressing \textbf{RAS}, \textbf{IO}, \textbf{Performance}, \textbf{Memory}, and \textbf{System Management}, trimming early bring-up failures by about 20\%.}

% Prior lab bullet duplicating OS names and validation pillars with diagnostic bullet (comment out; do not delete):
% \resumeItem{Deployed and maintained \textbf{Linux} and \textbf{Windows} GPU lab stacks spanning \textbf{BMC}, \textbf{BIOS}, and OS images plus repeatable workload staging for \textbf{system integration} bring-up, stressing \textbf{RAS}, \textbf{IO}, \textbf{Performance}, \textbf{Memory}, and \textbf{System Management}, trimming early bring-up failures by about 20\%.}

% Prior lab bullet with golden-image and heterogeneous phrasing (comment out; do not delete):
% \resumeItem{Deployed and maintained datacenter GPU rack fleets spanning \textbf{BMC}, \textbf{BIOS}, and heterogeneous OS image lines with repeatable workload staging ahead of \textbf{system integration} campaigns, tightening golden-image and provisioning consistency and trimming early bring-up failures by about 20\%.}

% Prior longer rack bring-up bullet (comment out; do not delete):
% \resumeItem{Deployed and maintained datacenter GPU test racks spanning \textbf{BMC} and \textbf{BIOS} updates, aligned OS images and pre-staged workloads before \textbf{system integration} runs, and cut early bring-up mistakes by about 20\%.}

% Prior rack bullet without UEFI flash narrative (comment out; do not delete):
% \resumeItem{Maintained datacenter GPU racks with \textbf{BMC} and \textbf{BIOS} updates, aligned OS images and staged workloads for \textbf{system integration}, cutting early bring-up failures by about 20\%.}

% Prior rack bullet with mixed-vendor phrasing and 20\% metric (comment out; do not delete):
% \resumeItem{Debugged mixed-vendor \textbf{BMC} setups and \textbf{UEFI}/\textbf{BIOS} mismatches on MI300X, MI325X, and MI350X racks, flashed and verified firmware images, and realigned OS lines so first-boot and image-mismatch failures during rack bring-up dropped by about 20\%.}

% Prior BMC/BIOS bullet without reuse metrics (comment out; do not delete):
% \resumeItem{Debugged \textbf{BMC} and \textbf{UEFI}/\textbf{BIOS} issues on MI300X, MI325X, and MI350X systems from different customer environments, flashed and verified firmware as needed, and wrote team-facing guides on BMC, BIOS, and flash bring-up that captured recurring pitfalls so others could repeat steps without rediscovering the same failures.}

% \resumeItem{Debugged \textbf{BMC} and \textbf{BIOS/UEFI} issues on MI300X, MI325X, and MI350X systems from different customer environments, flashed and verified firmware as needed, and wrote team-facing BMC, BIOS/UEFI, and flash bring-up guides teammates reused on \textbf{20+ rack bring-ups}, \textbf{cutting repeat setup clarifications by about 40\%} once pitfalls were centralized.}
% \resumeItem{Debugged \textbf{BMC} and \textbf{BIOS/UEFI} issues across MI300X, MI325X, and MI350X racks; documented flash and bring-up steps reused on \textbf{20+ rack bring-ups}, \textbf{cutting repeat setup clarifications by about 40\%}.}
% Prior customer configs phrasing (comment out; do not delete):
% \resumeItem{Debugged \textbf{BMC} and \textbf{BIOS/UEFI} failures on MI300X, MI325X, and MI350X racks from mixed customer configs; wrote flash and rack guides reused on \textbf{20+} bring-ups, \textbf{cutting repeat setup questions by about 40\%}.}
\resumeItem{Diagnosed \textbf{BMC} and \textbf{BIOS/UEFI} failures across hardware platforms, resolved complex firmware and boot issues, and documented standardized \textbf{hardware troubleshooting procedures} adopted across the team on \textbf{20+ platform bring-ups}; improved onboarding efficiency and reduced ambiguity for new hardware integration cycles.}

% Prior OOB bullet with tightening and inventory checks phrasing (comment out; do not delete):
% \resumeItem{Validated rack-level \textbf{out-of-band management} services for MI300X, MI325X, and MI350X systems next to \textbf{BMC}-driven diagnostics, tightening remote recovery and inventory checks so distributed lab racks stayed ready for recurring integration campaigns.}

% Prior longer OOB bullet (comment out; do not delete):
% \resumeItem{Validated rack-level \textbf{out-of-band management} on MI300X, MI325X, and MI350X racks alongside \textbf{BMC} health checks, keeping remote recovery and hardware inventory dependable for ongoing \textbf{system integration} work across the lab.}

% Prior short OOB validation bullet (comment out; do not delete):
% \resumeItem{Validated rack \textbf{out-of-band management} and \textbf{BMC} health on MI300X, MI325X, and MI350X for stable remote recovery, inventory, and \textbf{system integration}.}

% Prior OOB Python bullet without time metric (comment out; do not delete):
% \resumeItem{Implemented \textbf{Python} automation against rack out-of-band interfaces for MI300X, MI325X, and MI350X with \textbf{BMC} workflows for remote power, boot, and inventory during integration bring-up.}

% Prior OOB Python bullet with rack time-saving metric (comment out; similar time metrics vs launcher bullet; do not delete):
% \resumeItem{Implemented \textbf{Python} automation against rack out-of-band interfaces for MI300X, MI325X, and MI350X with \textbf{BMC} workflows for remote power, boot, and inventory during integration bring-up, \textbf{replacing about 25 minutes of manual console work per rack with under 5 minutes} of scripted actions for the routines I automated.}

% Prior launcher bullet without Git/Docker emphasis (comment out; do not delete):
% \resumeItem{Built \textbf{Python} tooling and an internal workload launcher that accelerated regression execution, cutting recurring manual setup from about 20 minutes to under 5 minutes for validation engineers.}

% Prior Python automation bullet naming Git (comment out; do not delete):
% \resumeItem{Extended \textbf{Python} test automation through an internal workload launcher and \textbf{Git}-tracked utilities so engineers reproduced diagnostics consistently, dropping recurring manual setup from about 20 minutes to under 5 minutes.}

% Prior Python automation bullet with version-controlled utilities phrasing (comment out; do not delete):
% \resumeItem{Extended \textbf{Python} test automation through an internal workload launcher and version-controlled utilities so engineers reproduced diagnostics consistently, dropping recurring manual setup from about 20 minutes to under 5 minutes.}

% Prior longer Python automation bullet (comment out; do not delete):
% \resumeItem{Extended \textbf{Python} automation and an internal workload launcher with shared, version-tracked helper scripts so teammates could rerun the same diagnostics reliably, cutting routine setup from about 20 minutes to under 5 minutes.}

% Prior launcher bullet without bolded time metric (comment out; do not delete):
% \resumeItem{Extended \textbf{Python} automation and an internal workload launcher with version-tracked scripts, shrinking diagnostic setup from about 20 minutes to under 5 minutes.}

% \resumeItem{Extended \textbf{Python} automation and an internal workload launcher with version-tracked scripts, shrinking diagnostic setup from \textbf{about 20 minutes to under 5 minutes} per run through reusable launch paths and AI-assisted edits in \textbf{Cursor} and \textbf{GitHub Copilot}.}
% \resumeItem{Extended \textbf{Python} automation and an internal workload launcher with \textbf{GitHub}-tracked scripts, shrinking diagnostic setup from \textbf{about 20 minutes to under 5 minutes} per run through reusable launch paths.}
% \resumeItem{Extended \textbf{Python} and \textbf{Bash} automation with \textbf{Git}-tracked diagnostic scripts and an internal workload launcher, shrinking recurring test-station setup from \textbf{about 20 minutes to under 5 minutes} per run through reusable execution paths.}
% \resumeItem{Built \textbf{Python} and \textbf{Bash} scripts with \textbf{Git}-tracked diagnostics and an internal workload launcher, cutting recurring test-station setup from \textbf{about 20 minutes to under 5 minutes} per run via reusable launch paths.}
% \resumeItem{Built \textbf{Python} and \textbf{Bash} scripts plus an internal workload launcher (\textbf{Git}-tracked) that cut manual diagnostic prep from \textbf{about 20 minutes to under 5 minutes} per run for the validation team.}
% Prior workload launcher phrasing (comment out; do not delete):
% \resumeItem{Built \textbf{Python} and \textbf{Bash} scripts plus an internal workload launcher that cut test launch time from \textbf{about 20 minutes to under 5 minutes} per run for the validation team.}
\resumeItem{Developed \textbf{Python} and \textbf{Bash}-based \textbf{test automation tooling} with internal launch frameworks and version-controlled procedures; influenced adoption across distributed validation teams, reducing manual test setup from \textbf{about 20 minutes to under 5 minutes} per cycle and improving consistency for new hardware platform bring-up.}

% Prior triage wording (comment out; do not delete):
% \resumeItem{Triaged hardware and software failures across ROCm, firmware, drivers, and interfaces with systematic log review and reproducible test flows, shortening time-to-root-cause on assigned issues by roughly 30\%.}

% Prior ROCm triage bullet without C/C++ (comment out; do not delete):
% \resumeItem{Triaged hardware and software defects spanning \textbf{ROCm}, firmware, drivers, and interconnect paths via systematic log review and scripted reproduction flows, cutting median time-to-root-cause on assigned defects by roughly 30\%.}

% Prior triage bullet naming ROCm (comment out; do not delete):
% \resumeItem{Triaged hardware and software defects spanning \textbf{ROCm}, \textbf{C/C++} firmware and driver stacks, and interconnect paths via systematic log review and scripted reproduction flows, cutting median time-to-root-cause on assigned defects by roughly 30\%.}

% Prior triage bullet with interconnect paths phrasing (comment out; do not delete):
% \resumeItem{Triaged hardware and software defects spanning \textbf{GPU} driver and firmware stacks, \textbf{C/C++} platform software, and interconnect paths via systematic log review and scripted reproduction flows, cutting median time-to-root-cause on assigned defects by roughly 30\%.}

% Prior longer triage bullet (comment out; do not delete):
% \resumeItem{Triaged hardware and software defects across \textbf{GPU} drivers, firmware, \textbf{C/C++} system code, and board-level links using logs and scripted reproduction steps, cutting median time to find root cause on assigned bugs by about 30\%.}

% Prior triage bullet without ROCm (comment out; do not delete):
% \resumeItem{Triaged defects across \textbf{GPU} drivers, firmware, \textbf{C/C++}, and interconnects using logs and repeatable test scripts, cutting median time-to-root-cause on assigned defects about 30\%.}

% Prior triage bullet with median time-to-root-cause metric (comment out; do not delete):
% \resumeItem{Triaged defects across \textbf{ROCm}, \textbf{GPU} drivers, firmware, \textbf{C/C++}, and interconnects using logs and repeatable test scripts, cutting median time-to-root-cause on assigned defects about 30\%.}

% Prior triage bullet with bold ROCm (comment out; do not delete):
% \resumeItem{Triaged defects across \textbf{ROCm}, \textbf{GPU} drivers, firmware, \textbf{C/C++}, and interconnects using logs; added \textbf{repeatable test scripts} so every handoff shipped the same repro bundle, which \textbf{cut reopened and duplicate defects on my queue by about 30\%} by eliminating one-off retests.}

% \resumeItem{Triaged defects across vendor \textbf{GPU} stacks including ROCm, firmware, \textbf{C/C++}, and interconnects using logs; added \textbf{repeatable test scripts} so every handoff shipped the same repro bundle, which \textbf{cut reopened and duplicate defects on my queue by about 30\%} by eliminating one-off retests.}
% \resumeItem{Triaged defects across \textbf{ROCm}, \textbf{GPU} driver/module stacks, firmware, \textbf{C/C++}, and interconnects using logs; added \textbf{repeatable test scripts} so each handoff shipped the same repro bundle, \textbf{cutting reopened and duplicate defects on my queue by about 30\%} by dropping one-off retests.}
% \resumeItem{Triaged defects across \textbf{ROCm}, \textbf{GPU} drivers, firmware, \textbf{C/C++}, and interconnects using logs and scripted repro steps; standardized handoff scripts with the same repro files, \textbf{cutting reopened defects on my queue by about 30\%}.}
% Prior ROCm-specific triage phrasing (comment out; do not delete):
% \resumeItem{Triaged defects across \textbf{ROCm} code and analyzed \textbf{C/C++} platform code for root cause; documented debug and repro steps so similar issues resolved faster, \textbf{cutting reopened defects on my queue by about 30\%}.}
\resumeItem{Triaged hardware and software integration failures across firmware, drivers, and platform code; developed and standardized reproducible \textbf{test procedures} and root-cause analysis documentation that influenced team practices, \textbf{reducing duplicate defect reruns by about 30\%} and accelerating resolution cycles for new platform validation.}
\resumeItem{Set up \textbf{virtualization environments} on datacenter GPU platforms using \textbf{KVM/QEMU}, validated guest workload execution under production-style loads, and documented setup procedures for team reuse across multiple platform generations.}
\resumeItem{Validated datacenter rack-level \textbf{out-of-band management} (BMC, remote power/boot, inventory services) for new hardware platforms alongside operational validation tests; coordinated with distributed teams to improve platform operational readiness and keep lab infrastructure reliable for recurring integration campaigns.}
% Prior GPU IP validation bullet (comment out; do not delete):
% \resumeItem{Ran \textbf{rocJPEG} workloads on MI systems and debugged \textbf{GFX} IP validation failures on the graphics and multimedia path during post-silicon passes.}

% Prior Confluence-forward cross-functional bullet (comment out; do not delete):
% \resumeItem{Partnered across cross-functional firmware, validation, and infrastructure teams using \textbf{Jira} for work tracking and \textbf{Confluence} for automation documentation and repeatable setup procedures, reducing misconfiguration and speeding adoption of lab automation across a distributed organization.}

% Prior wording before Confluence emphasis (comment out; do not delete):
% \resumeItem{Partnered with firmware and validation engineers across \textbf{Jira} and shared setup guides so teammates could repeat complex rack-style platform bring-up reliably in a distributed team environment.}

% Prior collaboration and documentation bullet (comment out; do not delete) — trimmed for one-page NVIDIA NCG layout:
% \resumeItem{Worked with infrastructure, firmware, and validation teams to track work in Jira and improve checklists and documentation, which made complex GPU diagnostic tests easier to reproduce and reduced misconfigurations in new test setups.}

% Original bullets (kept for reference)
% \resumeItem{Developed an internal \textbf{Python automation tool} that reduced a 20-minute daily validation and debug workflow to \textbf{5 minutes} for 4 engineers, saving \textbf{250+ hours annually}.}
% \resumeItem{Configured and supported \textbf{Linux-based datacenter GPU platforms}, installing workloads, validation tools, and system utilities to enable large-scale validation across distributed test environments.}
% \resumeItem{Performed \textbf{system-level debugging and failure triage} across Windows and Linux environments, investigating issues related to \textbf{RAS events, system resets, system management, workloads, and performance} to support root-cause analysis of GPU infrastructure platforms.}
% \resumeItem{Applied learnings from \textbf{system-level debug investigations} to improve \textbf{AI/HPC workload validation pipelines} for MI300X, MI325X, and MI350X GPUs, strengthening test coverage and accelerating issue debug and root-cause identification.}
% \resumeItem{Authored standardized debug and reproduction documentation and collaborated with infrastructure, \textbf{firmware}, and \textbf{validation} teams to resolve complex system issues across hardware, firmware, and software layers.}

% Prior Validation Engineer / GPU validation wording (comment out; do not delete):
% \resumeItem{Worked on post-silicon datacenter GPU system integration for MI300X, MI325X, and MI350X, running automated integration tests for performance and quality across 100+ test runs per week.}
% \resumeItem{Set up and ran tests on \textbf{Linux} and Windows GPU platforms with different BMC, BIOS, and OS images to check RAS, resets, IO, performance, and system management, cutting bring-up issues by about 20\%.}
% \resumeItem{Wrote \textbf{Python} scripts and a workload launch tool using an internal framework to speed up how quickly engineers could run regression tests, reducing manual setup time from about 20 minutes to under 5 minutes.}
% \resumeItem{Helped debug post-silicon issues involving ROCm, power-management firmware, drivers, and hardware by reading logs, reproducing problems, and sharing clear findings with senior engineers, shortening time-to-root-cause on assigned issues by roughly 30\%.}
% Prior variant without explicit \textbf{Jira} (comment out; do not delete):
% \resumeItem{Worked with infrastructure, firmware, and validation teams to track work in Jira and improve checklists and documentation, which made complex GPU tests easier to reproduce and reduced misconfigurations in new test setups.}

\resumeItemListEnd